\section{Related Work}
\label{sec:related_work}

\paragraph{Measurement error in causal inference.}
The epidemiology literature has long studied misclassification bias through causal diagrams~\cite{gustafson2004}. \citet{hernan_cole2009} formalized the starred-node DAG framework for reasoning about measurement bias, distinguishing differential from non-differential error. For binary variables, \citet{ogburn_vanderweele2012} proved that non-differential confounder misclassification produces adjusted estimates between the crude and true values, precluding sign reversal. \citet{brenner1993} showed this guarantee breaks for polytomous confounders, and \citet{veirod_laake2001} derived bias expressions for multi-category exposure misclassification in Poisson regression. \citet{edwards_lessler_naimi2021} further demonstrated that dependent misclassification of covariate and exposure can bias in either direction. Our work extends this line by providing a systematic taxonomy across seven DAG topologies with closed-form expressions for the $K$-class case, identifying which topologies preserve and which break the sign-flip guarantee.

\paragraph{Text classifier errors in causal studies.}
\citet{wood_doughty2018} first applied DAG reasoning to text classifier measurement error, noting matrix adjustment correction but without formal bias theorems. \citet{teblunthuis2024} provided the most complete binary treatment: four DAG scenarios with testable conditions and a Maximum Likelihood Adjustment method, but restricted to binary classification. \citet{yang_adomavicius2018} formalized the two-stage pipeline problem (data mining followed by econometric analysis) and proved that ML-generated variables systematically bias downstream regression. \citet{keith_jensen_oconnor2020} surveyed the text-and-causal-inference landscape, identifying classifier measurement error as a key open challenge. \citet{lee_wood_doughty2024} extended to categorical LLM classifiers on clinical data but covered only one topology (unobserved confounding). \citet{chen_bhattacharya_keith2024} used zero-shot NLP models as proxies in proximal causal inference, proving identification conditions for one DAG structure. We generalize across all seven topologies with multi-class confusion matrices, providing both bias expressions and sensitivity analysis tools absent from prior work.

\paragraph{LLM annotation quality.}
Recent work has quantified the downstream risks of LLM annotation. \citet{baumann2025} demonstrated across 37 tasks that 31\% of hypotheses yield incorrect conclusions with LLM annotations. \citet{kim_garg2025} showed that LLM errors are 60\% correlated across models, undermining ensemble-based mitigation. \citet{messing2026} decomposed LLM pipeline variance via linear mixed models but explicitly excluded causal identification. \citet{stolwijk_boukes_trilling2025} found that LLM annotations differ from human annotations in prevalence even when per-sample accuracy is adequate, and \citet{cheng_mayya_sedoc2024} decomposed annotation measurement error into four source types. These works diagnose the annotation quality problem; our contribution provides the causal-analytic bridge connecting annotation error structure to downstream causal estimate validity.

\paragraph{Sensitivity analysis for observational studies.}
Our ASV framework draws on the tradition of sensitivity analysis for unmeasured confounding~\cite{rosenbaum2002}. \citet{vanderweele_ding2017} introduced E-values as interpretable sensitivity measures, and \citet{ding_vanderweele2016} developed assumption-free sensitivity bounds. \citet{vanderweele_li2019} extended sensitivity analysis to differential measurement error but only for binary settings. \citet{smith_vanderweele2021} proposed bounds for multiple bias sources including misclassification. Our ASVs adapt the E-value philosophy (``how much bias would be needed to explain away the finding?'') to annotation-induced bias, with three complementary measures tailored to different topology-dependent failure modes.

\paragraph{Correction methods.}
Prediction-powered inference~\cite{angelopoulos2023} and design-based supervised learning~\cite{egami_hinck_stewart_wei2023, gligoric2025} provide general frameworks for valid inference with ML-generated labels, combining predicted labels with small gold-standard samples. \citet{zrnic_candes2024} extended PPI via cross-prediction for improved efficiency. \citet{audinet2025} benchmarked these methods across datasets. These correction methods are complementary to our framework: our topology-dependent ASV analysis identifies \emph{when} correction is needed and \emph{which} failure mode to address, while PPI/DSL provide the \emph{how}.
